% Page 052 — page imprimée 49
% Contenu seul : le préambule, l'en-tête et la pagination sont gérés par meyrueis.tex

\textbf{\Large L’ordre moral}

Le général commandant les chantiers, catholique pratiquant, voua au Maréchal
Pétain un attachement inconditionnel qu’il exigeait de ses subordonnés.

\emph{« Le père Forestier, partisan de l’obéissance à Vichy et plus tard du départ des
jeunes des chantiers en Allemagne pour Le Service du Travail Obligatoire (STO) est
nommé aumônier général, tout en restant celui des Scouts de France. Il voit en
Pétain une grande chance pour le catholicisme français... tout comme son chef le
Général de la porte du Theil qui s’exprime ainsi :}

\begin{quote}
\emph{« Il n’y a de sûrs que les ordres des chefs régulièrement investis d’un mandat et fidèlement attachés
par la voie hiérarchique au maréchal, seul responsable et seul garant, ne l’oublions jamais, de l’unité
de la Patrie, parce que c’est le maréchal et parce qu’il a reçu mission de nous conduire, donc que les
lumières pour cela ne lui manqueront jamais ». (15 Nov. 1942)}
\end{quote}

Il est intéressant de relever quelques décisions du Conseil Municipal de Meyrueis à la même époque :

\begin{description}
\item[9 Février 1941 :] Désignation de la place du Mal. Pétain au lieu de la place d’Orléans. (unanimité)
\item[6 Décembre 1942 :] Inébranlable confiance et absolue fidélité au Mal. Pétain (à l’unanimité)
\item[9 Août 1942 :] Augmentation du crédit de gardiennage de l’Eglise qui passe de 200F à 1000F.
\item[25 Juillet 1943 :] Reconnaissance légale de la Congrégation des Frères de la Doctrine Chrétienne.
\item[19 Décembre 43 :] Décision d’abattre le dernier orme de Sully…
\end{description}

\emph{De l’amalgame entre la religion et le Vichysme naît une morale de confection où
l’on retrouve tous les poncifs de la Révolution Nationale : culte du Maréchal,
expiation masochiste des fautes de la République, retour à la vivifiante nature,
résurrection du corporatisme. Furent nommés 4 ou 5 aumôniers protestants dans
toute la région et 20 à 30 aumôniers catholiques. Ces derniers sont partout. Ils
jouent à l’occasion le rôle d’éminence grise auprès des commissaires. Certains, bien
sûr, sont de haut niveau : on dit la messe souvent et partout, en plein air comme
dans les foyers. Jean Delage donne parfaitement le ton de ces « messes chantiers » :
« Le drapeau salué, le clairon sonne, les garçons s’agenouillent : c’est l’élévation ».
Certes nul n’est tenu de suivre ces manifestations religieuses ; mais il est difficile de
les éviter ». (L’espelido : Histoire des Chantiers de la Jeunesse en Languedoc-Roussillon : Pierre
Mazier p 83 et suivantes…)}

Certains résistèrent à leur façon en organisant un bal du dimanche soir au Planet dans l’ancien four
banal. Autour d’un vieux phono se retrouvaient garçons et filles du village et.. de la vallée du
Béthuzon.. (Tous les bals avaient été interdits après l’armistice)

Mais quelques uns ne partagent pas les vues de l’aumônier général, ainsi :

\emph{« Georges Vidal, assistant élève à l’Ecole Régionale des Cadres , naïvement, le lui dit. Presque
aussitôt son stage est interrompu et il est renvoyé sans ménagement dans son camp perdu au-dessus
de Lodève ; motif « ne peut faire un chef de groupe »... (P. Mazier op. cité)}

Pour le groupement 19 avaient été nommés trois aumôniers catholiques Navarre,
Beys , Durand et un aumônier protestant Georges Roth. Ce dernier, alsacien, avait été
pasteur au Havre, officier d’artillerie en 1939-1940. Il fut dénoncé par ses supérieurs
pour avoir contesté le devoir d’obéissance au gouvernement de Vichy. (voir lettre jointe)
